% !TeX root = ../../Book2.tex
\section{能量估计与低正则迭代}
\label{b2:sec:H201}

系数只有可测性时，不能把方程逐项微分，却仍可以改变测试函数。截去低于某个高度的部分，可以观察高值集合；提高解的幂次，可以逐步增强可积性；取倒数或对数，则能接近正解的下界。这些操作都发生在弱形式的一阶导数之内。本节先把测试的合法性说清楚，再沿正幂、对数和倒数三条线建立线性 Harnack 比较。

\ProseParagraphBreak

以下讨论实标量方程。设 \(d\geq2\)，\(\Omega\subset\mathbb R^d\) 开，实矩阵场 \(A:\Omega\to\mathbb R^{d\times d}\) 可测，并满足
\[
 A(x)\xi\cdot\xi\geq\lambda|\xi|^2,
 \quad |A(x)\xi|\leq\Lambda|\xi|
 \quad(\xi\in\mathbb R^d)
\]
于几乎处处的 \(x\)，其中 \(0<\lambda\leq\Lambda<\infty\)。不要求 \(A\) 对称，也不使用系数的导数。实函数 \(u\in H^1_{\mathrm{loc}}(\Omega)\) 是齐次方程的弱次解，是指
\[
 \int_\Omega A\nabla u\cdot\nabla\varphi\leq0
 \quad(0\leq\varphi\in C_c^\infty(\Omega));
\]
若对全部实值 \(\varphi\in C_c^\infty(\Omega)\) 等式成立，则称为弱解。这个符号约定对应 \(-\operatorname{div}(A\nabla u)\leq0\)。下文的上、下界均先按几乎处处的意义理解，直到最后使用连续代表。

\subsection{Caccioppoli 不等式}
\label{b2:subsec:H20101}

把解本身或它的正部放到测试位置时，所得函数通常不光滑。弱形式可以接受这种测试，但理由来自 Sobolev 逼近以及通量的平方可积性，不能只靠形式上的分部积分。

\Needspace{6\baselineskip}
\begin{lemma}[非负 Sobolev 测试与正部]
\label{b2:lem:H2-sobolev-tests}
设 \(d\geq2\)，\(\Omega\subset\mathbb R^d\) 开，\(A:\Omega\to\mathbb R^{d\times d}\) 可测，且对几乎处处的 \(x\) 和全部 \(\xi\in\mathbb R^d\)，有
\[
 A(x)\xi\cdot\xi\geq\lambda|\xi|^2,
 \quad |A(x)\xi|\leq\Lambda|\xi|,
 \quad 0<\lambda\leq\Lambda<\infty.
\]
若实函数 \(u\in H^1_{\mathrm{loc}}(\Omega)\) 满足
\(\int_\Omega A\nabla u\cdot\nabla\varphi\leq0\) 对所有非负 \(\varphi\in C_c^\infty(\Omega)\) 成立，则同一不等式对每个非负、具有紧支集的 \(\varphi\in H^1(\Omega)\) 成立。对任意 \(k\in\mathbb R\)，函数 \((u-k)_+\) 也满足这一弱次解不等式。
\end{lemma}
\begin{proof}
先设非负 \(\varphi\in H^1(\Omega)\) 的支集包含于 \(K\Subset\Omega\)。在 \(K\) 附近取等于一的光滑紧支集函数，将 \(\varphi\) 乘以它再零延拓，便得到 \(H^1(\mathbb R^d)\) 函数。用非负磨光核卷积，得到 \(\varphi_\delta\geq0\)。当 \(\delta\) 足够小时，这些光滑函数的支集都落在同一个 \(U\Subset\Omega\) 内，并且 \(\varphi_\delta\to\varphi\) 于 \(H^1(U)\)。由于
\[
 \left|\int_U A\nabla u\cdot\nabla(\varphi_\delta-\varphi)\right|
 \leq\Lambda\|\nabla u\|_{L^2(U)}
       \cdot\|\nabla(\varphi_\delta-\varphi)\|_{L^2(U)},
\]
弱次解不等式可取极限。若原来是弱解，同一逼近不需要保留符号，因而允许任意实值的紧支撑 \(H^1\) 测试。

再令 \(v=(u-k)_+\)。Sobolev 截断链式法则给出
\(\nabla v=\mathbf1_{\{u>k\}}\cdot\nabla u\)。固定非负 \(\psi\in C_c^\infty(\Omega)\)，取分段线性函数
\[
 \theta_\delta(t)=
 \begin{cases}
 0,&t\leq0,\\
 t/\delta,&0<t<\delta,\\
 1,&t\geq\delta.
 \end{cases}
\]
函数 \(\theta_\delta(u-k)\cdot\psi\) 是合法的非负 Sobolev 测试。代入后得到
\[
 \int\theta_\delta(u-k)A\nabla u\cdot\nabla\psi
 +\int\psi\cdot\theta_\delta'(u-k)
                      A\nabla u\cdot\nabla u\leq0.
\]
第二项非负，故可从左边删除。第一项由紧支集上的
\(\Lambda|\nabla u|\cdot|\nabla\psi|\) 控制，这是可积函数；令 \(\delta\downarrow0\)，得到
\(\int A\nabla v\cdot\nabla\psi\leq0\)。分段函数折点上的导数取值不影响结果：Sobolev 函数在每个固定水平集上的梯度几乎处处为零，这也是截断链式法则的一部分。
\end{proof}

正部保留次解性质，使高度截断可以反复使用。下面的估计把弱形式中的椭圆性变成梯度控制；截断函数的导数记录了局部化所付的代价。

\Needspace{6\baselineskip}
\begin{proposition}[水平截断的 Caccioppoli 估计]
\label{b2:prop:H2-level-caccioppoli}
设 \(d\geq2\)，\(\Omega\subset\mathbb R^d\) 开，\(A:\Omega\to\mathbb R^{d\times d}\) 为可测实矩阵场，且
\(A(x)\xi\cdot\xi\geq\lambda|\xi|^2\)、\(|A(x)\xi|\leq\Lambda|\xi|\) 对几乎处处的 \(x\) 及全部 \(\xi\in\mathbb R^d\) 成立，其中 \(0<\lambda\leq\Lambda<\infty\)。设实函数 \(u\in H^1_{\mathrm{loc}}(\Omega)\) 满足
\(\int A\nabla u\cdot\nabla\varphi\leq0\) 对每个非负 \(\varphi\in C_c^\infty(\Omega)\) 成立。对 \(k\in\mathbb R\)、\(w=(u-k)_+\) 及 \(0\leq\eta\leq1\)、\(\eta\in C_c^\infty(\Omega)\)，有
\[
 \int_\Omega\eta^2|\nabla w|^2
 \leq4\frac{\Lambda^2}{\lambda^2}
          \int_\Omega w^2|\nabla\eta|^2
\]
以及
\begin{equation}\label{b2:eq:H2-level-energy}
 \int_\Omega|\nabla(\eta\cdot w)|^2
 \leq C\int_\Omega w^2|\nabla\eta|^2,
\end{equation}
其中 \(C\) 只依赖 \(\lambda,\Lambda\)。特别地，若
\(0<r<s\)、\(B_s=B(x_0,s)\Subset\Omega\)，则
\[
 \int_{B_r}|\nabla(u-k)_+|^2
 \leq\frac{C}{(s-r)^2}\int_{B_s}(u-k)_+^2.
\]
\end{proposition}
\begin{proof}
由上一引理，可用 \(\eta^2\cdot w\) 测试。在 \(w>0\) 处有 \(\nabla w=\nabla u\)，在其余部分 \(w=0\)、\(\nabla w=0\) 几乎处处，所以
\[
 \lambda\int\eta^2|\nabla w|^2
 \leq2\Lambda\int\eta w|\nabla w|\cdot|\nabla\eta|
 \leq\frac\lambda2\int\eta^2|\nabla w|^2
       +\frac{2\Lambda^2}{\lambda}\int w^2|\nabla\eta|^2.
\]
把梯度项移回左边即得第一式。再由
\[
 |\nabla(\eta\cdot w)|^2
 \leq2\eta^2|\nabla w|^2+2w^2|\nabla\eta|^2
\]
得到 \eqref{b2:eq:H2-level-energy}。最后取 \(\eta=1\) 于 \(B_r\)，支撑于 \(B_s\)，并使 \(|\nabla\eta|\leq C_d/(s-r)\)，即得两球估计。
\end{proof}

对齐次弱解，\(-u\) 也是解，因而负部可按相同方式估计；若改用 \(\eta^2\cdot(u-c)\) 测试，还可估计整个梯度。这一能量界不含 \(\nabla A\)，而\chapref{b2:ch:16}中的差商论证需要额外的系数正则性。两者控制的对象不同：这里从函数的局部大小控制一阶能量，差商方法则试图产生二阶弱导数。

\ProseParagraphBreak

齐次假设也有实际作用。若 \(-\operatorname{div}(A\nabla u)=f\)，测试会多出 \(\int f\eta^2w\)。必须先说明该积分有意义，再估计它；只有这样才能把源项的可积指数和缩放带入迭代。将这一项删除，会把含源项方程误当成齐次方程。

\subsection{De Giorgi 迭代}
\label{b2:subsec:H20102}

\cref{b2:thm:16-local-boundedness}已经对有界可测一致椭圆系数的实弱次解，完整证明了
\[
 \operatorname*{ess\,sup}_{B_R}u_+
 \leq C\left(R^{-d}\int_{B_{2R}}u_+^2\right)^{1/2}.
\]
该证明只使用水平截断能量、Sobolev 不等式及一个数列迭代，不需要 Harnack 比较。即使解只属于局部 \(H^1\)，也可先在稍小的外球上应用结论，再把外半径增大；若右侧有限，单调性给出所列形式。

\ProseParagraphBreak

这一过程称为 \term[DeGiorgiDiedai]{De Giorgi 迭代}[De Giorgi iteration]。它同时收缩半径、提高截断高度。若第 \(j\) 轮的高度为 \(k_j\)，则下一轮高值集合上有
\((u-k_j)_+\geq k_{j+1}-k_j\)，所以旧的能量还能控制新集合的测度。把这一测度界与 Sobolev 的积分增益结合，归一化后的能量满足
\[
 Y_{j+1}\leq Cb^jY_j^{1+\alpha},
 \quad b>1,\quad\alpha>0.
\]
起始高度由初始平方平均选择，使 \(Y_0\) 足够小，才会得到 \(Y_j\to0\)。超线性幂 \(1+\alpha\) 是小量改善的来源；若只有 \(Y_{j+1}\leq Cb^jY_j\)，这条论证不能迫使高值部分消失。

\ProseParagraphBreak

本节不重复这条已完成的高度迭代。下面改用幂函数，以便进一步接入倒数和对数测试。两种路线的区别也解释了方法的选择：障碍问题天然允许单侧高度截断；正解比较则需要把函数的大值与小值联系起来，幂函数的正、负方向更便于配合。

\subsection{Moser 迭代}
\label{b2:subsec:H20103}

逐次提高可积指数的路线称为 \term[MoserDiedai]{Moser 迭代}[Moser iteration]。形式上，用 \(u^{q-1}\) 测试会出现 \(u^{q-2}|\nabla u|^2\)，它正是 \(u^{q/2}\) 的梯度平方乘以一个常数。但本轮只知道 \(u\in L^q\)，并不知道这个幂已属于 \(H^1\)。因此必须先把幂函数截成全局 Lipschitz 函数，取得独立于截断高度的估计，再放开截断。

\ProseParagraphBreak

先记下本节使用的 Sobolev 指数：当 \(d>2\) 时取 \(\chi=d/(d-2)\)，当 \(d=2\) 时取 \(\chi=2\)。对 \(z\in H_0^1(B(0,2))\)，有
\(\|z\|_{2\chi}^2\leq C_d\|\nabla z\|_2^2\)。二维情形可先用零延拓后的 \(W^{1,4/3}\) 到 \(L^4\) 的 Sobolev 不等式，再用有限体积上的 Hölder 不等式把梯度的 \(L^{4/3}\) 范数界为 \(L^2\) 范数。缩放回任意 \(B_{2R}=B(x_0,2R)\)，所得形式是
\begin{equation}\label{b2:eq:H2-scaled-sobolev}
 \left(R^{-d}\int_{B_{2R}}|z|^{2\chi}\right)^{1/\chi}
 \leq C_dR^{2-d}\int_{B_{2R}}|\nabla z|^2
 \quad(z\in H_0^1(B_{2R})).
\end{equation}
这一个式子包含了两个维数范围。特别是在二维，若改写成未归一化的范数，右侧会带一个 \(R\)；不能把单位球常数原样用于任意半径。

\Needspace{6\baselineskip}
\begin{lemma}[幂函数能量估计]
\label{b2:lem:H2-power-energy}
设 \(d\geq2\)，\(x_0\in\mathbb R^d\)、\(R>0\)，记 \(B_t=B(x_0,t)\)。可测实矩阵场 \(A\) 在 \(B_{2R}\) 中满足
\(A(x)\xi\cdot\xi\geq\lambda|\xi|^2\)、\(|A(x)\xi|\leq\Lambda|\xi|\)，其中 \(0<\lambda\leq\Lambda<\infty\)，两式对几乎处处的 \(x\) 和全部 \(\xi\in\mathbb R^d\) 成立。设实函数 \(u\in H^1_{\mathrm{loc}}(B_{2R})\) 几乎处处非负，且
\(\int A\nabla u\cdot\nabla\varphi\leq0\) 对所有非负 \(\varphi\in C_c^\infty(B_{2R})\) 成立。取 \(q\geq2\)、\(R\leq r<s\leq2R\)，并假定 \(u\in L^q(B_s)\)。令 \(\chi=d/(d-2)\) 当 \(d>2\)，\(\chi=2\) 当 \(d=2\)。则
\[
 \left(R^{-d}\int_{B_r}u^{q\chi}\right)^{1/(q\chi)}
 \leq
 \left(\frac{Cq^2R^2}{(s-r)^2}\right)^{1/q}
 \left(R^{-d}\int_{B_s}u^q\right)^{1/q},
\]
其中 \(C\) 只依赖 \(d,\lambda,\Lambda\)。
\end{lemma}
\begin{proof}
固定 \(L>0\)。当 \(q>2\) 时，对 \(t\geq0\) 定义
\[
 h_L(t)=t\min(t,L)^{q-2},
 \quad w_L(t)=t\min(t,L)^{(q-2)/2};
\]
当 \(q=2\) 时直接定义 \(h_L(t)=w_L(t)=t\)。这样零点处没有 \(0^0\) 的约定问题。两函数在 \([0,\infty)\) 上 Lipschitz，零点取零，并可延拓为实轴上的 Lipschitz 函数。对于 \(q>2\)，其两支导数为
\[
 \begin{array}{c|cc}
 &0<t<L&t>L\\ \hline
 h_L'(t)&(q-1)t^{q-2}&L^{q-2}\\
 w_L'(t)&\dfrac q2t^{(q-2)/2}&L^{(q-2)/2}.
 \end{array}
\]
由这两行直接得到
\[
 h_L(t)^2\leq w_L(t)^2\cdot h_L'(t),
 \quad |w_L'(t)|^2\leq q^2\cdot h_L'(t),
 \quad h_L'(t)\geq0.
\]
第一式采用不含除法的形式，因而在 \(t=0\) 时仍有意义；\(q=2\) 时三式同样成立。折点 \(t=L\) 的导数值任取一侧即可，因为 \(\nabla u=0\) 几乎处处于 \(\{u=L\}\)。

取 \(0\leq\eta\leq1\)、\(\eta\in C_c^\infty(B_s)\)，且 \(\eta=1\) 于 \(B_r\)、\(|\nabla\eta|\leq C_d/(s-r)\)。对固定 \(L\)，函数 \(\eta^2\cdot h_L(u)\) 是合法的非负 Sobolev 测试。弱次解不等式给出
\[
 \begin{aligned}
 \lambda\int\eta^2h_L'(u)|\nabla u|^2
 &\leq2\Lambda\int\eta\cdot h_L(u)
                        |\nabla u|\cdot|\nabla\eta|\\
 &\leq\frac\lambda2\int\eta^2h_L'(u)|\nabla u|^2
       +\frac{2\Lambda^2}{\lambda}
                    \int w_L(u)^2|\nabla\eta|^2.
 \end{aligned}
\]
这里使用了 \(h_L\leq w_L\sqrt{h_L'}\)，再对两个非负因子使用 Young 不等式。吸收左边的一半，并用 \(|w_L'|^2\leq q^2h_L'\)，得到
\[
 \int|\nabla(\eta\cdot w_L(u))|^2
 \leq Cq^2\int w_L(u)^2|\nabla\eta|^2
 \leq\frac{Cq^2}{(s-r)^2}\int_{B_s}u^q.
\]
最后一步来自 \(w_L(u)^2\leq u^q\)。对 \(\eta\cdot w_L(u)\in H_0^1(B_{2R})\) 使用 \eqref{b2:eq:H2-scaled-sobolev}，得到
\[
 \left(R^{-d}\int_{B_r}w_L(u)^{2\chi}\right)^{1/\chi}
 \leq\frac{Cq^2R^2}{(s-r)^2}
              R^{-d}\int_{B_s}u^q.
\]
当 \(L\to\infty\) 时，\(w_L(u)\to u^{q/2}\) 几乎处处。左边用 Fatou 引理，右边保持不变，再取 \(q\) 次根即得结论。整个过程只使用本轮已有的 \(L^q\) 可积性，没有预设下一轮幂函数的梯度可积。
\end{proof}

\Needspace{6\baselineskip}
\begin{theorem}[正幂迭代的局部上界]
\label{b2:thm:H2-moser-upper}
设 \(d\geq2\)，\(x_0\in\mathbb R^d\)、\(R>0\)，记 \(B_t=B(x_0,t)\)。可测实矩阵场 \(A\) 在 \(B_{2R}\) 上满足
\(A(x)\xi\cdot\xi\geq\lambda|\xi|^2\)、\(|A(x)\xi|\leq\Lambda|\xi|\)，其中 \(0<\lambda\leq\Lambda<\infty\)，两式对几乎处处的 \(x\) 和全部实 \(\xi\) 成立。设实函数 \(u\in H^1_{\mathrm{loc}}(B_{2R})\cap L^2(B_{2R})\) 几乎处处非负，并满足
\(\int A\nabla u\cdot\nabla\varphi\leq0\) 对所有非负 \(\varphi\in C_c^\infty(B_{2R})\) 成立。则
\[
 \operatorname*{ess\,sup}_{B_R}u
 \leq C\left(R^{-d}\int_{B_{2R}}u^2\right)^{1/2},
\]
其中 \(C\) 只依赖 \(d,\lambda,\Lambda\)。
\end{theorem}
\begin{proof}
置 \(v(y)=u(x_0+Ry)\)、\(\widetilde A(y)=A(x_0+Ry)\)。在弱形式中换元，出现的共同正因子 \(R^{d-2}\) 可以约去；\(\widetilde A\) 保留同一组椭圆界。因此先证 \(R=1\)、\(x_0=0\) 的结论即可。

令 \(\chi\) 如幂函数能量估计所定，并取
\[
 q_j=2\chi^j,
 \quad r_j=1+2^{-j}
 \quad(j=0,1,2,\ldots).
\]
从 \(q_0=2\)、\(r_0=2\) 开始，上一引理逐次保证下一轮所需的可积性，而且
\[
 \|u\|_{L^{q_{j+1}}(B_{r_{j+1}})}
 \leq\left(Cq_j^2\cdot2^{2j+2}\right)^{1/q_j}
          \|u\|_{L^{q_j}(B_{r_j})}.
\]
把各轮相乘。因为
\[
 \sum_{j=0}^\infty
 \frac{\log C+2\log q_j+(2j+2)\log2}{q_j}<\infty,
\]
乘积被一个只依赖 \(d,\lambda,\Lambda\) 的常数控制。于是
\(\|u\|_{L^{q_j}(B_1)}\leq K\)，其中 \(K=C\|u\|_{L^2(B_2)}\)，与 \(j\) 无关。若存在 \(t>K\) 使 \(E=\{x\in B_1:u(x)>t\}\) 具有正测度，则
\[
 K\geq\|u\|_{L^{q_j}(B_1)}\geq t|E|^{1/q_j}
 \longrightarrow t,
\]
产生矛盾。因此 \(u\leq K\) 几乎处处。换元回原来的球，
\(\|v\|_{L^2(B_2)}=R^{-d/2}\|u\|_{L^2(B_{2R})}\)，即得所列估计。
\end{proof}

对有符号弱次解，可先由\cref{b2:lem:H2-sobolev-tests}取正部；对弱解再分别处理 \(u_+\) 和 \((-u)_+\)，便得到绝对值的局部上界。到这里，幂迭代已到达与 De Giorgi 迭代相同的结论。下一步需要把平方平均换成很小的正幂平均，因为对数的指数可积性只会产生一个小指数。

\Needspace{6\baselineskip}
\begin{lemma}[局部上界的低次平均版本]
\label{b2:lem:H2-small-power}
设 \(d\geq2\)，\(x_0\in\mathbb R^d\)、\(R>0\)，记 \(B_t=B(x_0,t)\)。可测实矩阵场 \(A\) 在 \(B_{2R}\) 中满足
\(A(x)\xi\cdot\xi\geq\lambda|\xi|^2\)、\(|A(x)\xi|\leq\Lambda|\xi|\)，其中 \(0<\lambda\leq\Lambda<\infty\)，两式对几乎处处的 \(x\) 和全部实 \(\xi\) 成立。实函数 \(u\in H^1_{\mathrm{loc}}(B_{2R})\) 几乎处处非负，并且
\(\int A\nabla u\cdot\nabla\varphi\leq0\) 对所有非负 \(\varphi\in C_c^\infty(B_{2R})\) 成立。对每个 \(a>0\)，有
\[
 \operatorname*{ess\,sup}_{B_R}u
 \leq C_a\left(R^{-d}\int_{B_{2R}}u^a\right)^{1/a},
\]
其中 \(C_a\) 只依赖 \(a,d,\lambda,\Lambda\)；右侧为无穷时按通常的扩展实数意义理解。
\end{lemma}
\begin{proof}
只需考虑右侧有限的情形，并缩放到 \(R=1\)。当 \(a\geq2\) 时，有限体积上的 Hölder 不等式保证 \(u\in L^2(B_2)\)，由上一项平方平均上界即得结论。

设 \(0<a<2\)，记 \(M(r)=\operatorname*{ess\,sup}_{B_r}u\)。首先，\(M(3/2)<\infty\)：用有限个较小球覆盖 \(\overline{B_{3/2}}\)，使每个球的二倍球都紧包含于 \(B_2\)，再逐球使用平方平均上界。所需平方积分来自 \(u\in H^1_{\mathrm{loc}}(B_2)\)。这个有限性先于后面的吸收步骤。

对 \(1\leq r<s\leq3/2\)，在 \(\overline{B_r}\) 各点放置半径 \(\rho=(s-r)/4\) 的球，其二倍球均在 \(B_s\) 内。取有限子覆盖，在每个球上使用已证估计，得
\[
 \begin{aligned}
 M(r)&\leq C(s-r)^{-d/2}\|u\|_{L^2(B_s)}\\
 &\leq C(s-r)^{-d/2}M(s)^{1-a/2}
               \left(\int_{B_{3/2}}u^a\right)^{1/2}.
 \end{aligned}
\]
若 \(M(s)=0\)，结论已成立。否则令
\(U=(\int_{B_{3/2}}u^a)^{1/a}\)，对上式用指数 \(2/(2-a)\) 与 \(2/a\) 的带参数 Young 不等式。即使 \(a<1\)，这个积分幂量仍有意义，且此处不使用范数三角不等式。对任意 \(0<\vartheta<1\)，得到
\[
 M(r)\leq\vartheta M(s)
       +C_{a,\vartheta}(s-r)^{-d/a}U.
\]
取 \(r_j=3/2-2^{-j-1}\)，于是 \(r_0=1\)，且
\(r_{j+1}-r_j=2^{-j-2}\)。反复代入可得
\[
 M(1)\leq\vartheta^nM(r_n)
   +C_{a,\vartheta}2^{2d/a}U
        \sum_{j=0}^{n-1}\left(\vartheta2^{d/a}\right)^j.
\]
选择 \(\vartheta<2^{-d/a}\)，级数可和；由于
\(M(r_n)\leq M(3/2)<\infty\)，第一项趋零。故
\(M(1)\leq C_aU\)，扩大右侧积分域到 \(B_2\)，再恢复尺度即得结论。这里没有把未知的无穷上确界移到不等式另一边。
\end{proof}

仅有正幂估计仍然不能给出正下界。对正解取对数后，加法常数对应原函数的乘法大小，而平均振荡消去了这个常数。因此，对数能量适合连接正幂与负幂：先控制与整体高度无关的振荡，再用指数可积性把它还原成两个方向的平均。

\Needspace{6\baselineskip}
\begin{proposition}[对数能量与正负平均连接]
\label{b2:prop:H2-log-bridge}
设 \(d\geq2\)，\(\Omega\subset\mathbb R^d\) 开，\(A:\Omega\to\mathbb R^{d\times d}\) 为可测实矩阵场，且
\(A(x)\xi\cdot\xi\geq\lambda|\xi|^2\)、\(|A(x)\xi|\leq\Lambda|\xi|\) 对几乎处处的 \(x\) 和全部实 \(\xi\) 成立，其中 \(0<\lambda\leq\Lambda<\infty\)。设实函数 \(u\in H^1_{\mathrm{loc}}(\Omega)\) 几乎处处非负，且
\(\int A\nabla u\cdot\nabla\varphi=0\) 对每个实值 \(\varphi\in C_c^\infty(\Omega)\) 成立。若 \(Q\) 为轴平行立方体，\(4Q\Subset\Omega\)，则存在只依赖 \(d,\lambda,\Lambda\) 的 \(a_0>0\)，使每个 \(\varepsilon>0\)、\(v=u+\varepsilon\) 满足
\[
 \left(\frac1{|Q|}\int_Qv^{a_0}\right)
 \left(\frac1{|Q|}\int_Qv^{-a_0}\right)\leq9.
\]
此外，对每个实值 \(\eta\in C_c^\infty(\Omega)\)，有
\[
 \int_\Omega\eta^2|\nabla\log v|^2
 \leq4\frac{\Lambda^2}{\lambda^2}
                  \int_\Omega|\nabla\eta|^2.
\]
这里 \(4Q\) 表示与 \(Q\) 同中心、边长为四倍的立方体。
\end{proposition}
\begin{proof}
常数平移不改变梯度，所以 \(v\) 仍为齐次弱解。因为 \(v\geq\varepsilon\)，倒数函数在其取值范围内及其一阶导数都有界，\(\eta^2/v\) 属于紧支撑的 \(H^1\)，可作为测试。又有 \(\log v\in H^1_{\mathrm{loc}}(\Omega)\)：负部由 \(\varepsilon\) 控制，正部可由 \(u\) 加一个依赖 \(\varepsilon\) 的常数控制，而链式法则给出 \(\nabla\log v=\nabla u/v\)。弱方程因此给出
\[
 \begin{aligned}
 \lambda\int\eta^2|\nabla\log v|^2
 &\leq2\Lambda\int|\eta|\cdot|\nabla\log v|
                                  \cdot|\nabla\eta|\\
 &\leq\frac\lambda2\int\eta^2|\nabla\log v|^2
                  +\frac{2\Lambda^2}{\lambda}
                             \int|\nabla\eta|^2.
 \end{aligned}
\]
吸收后得到对数能量界，常数与 \(\varepsilon\) 无关。

令 \(b=\log v\)。对 \(Q\) 的任意轴平行子立方体 \(P\)，包括 \(P=Q\)，取在 \(P\) 上为一、支撑于 \(2P\) 的截断，且梯度不超过 \(C_d/\ell(P)\)。所有 \(2P\) 都包含于 \(3Q\)，所以这些测试统一合法。能量界给出
\[
 \int_P|\nabla b|^2\leq C\ell(P)^{d-2}.
\]
立方体的 Poincaré 不等式和 Cauchy--Schwarz 不等式于是给出
\[
 \begin{aligned}
 \frac1{|P|}\int_P|b-b_P|
 &\leq\left(\frac1{|P|}\int_P|b-b_P|^2\right)^{1/2}\\
 &\leq C_d\ell(P)
          \left(\frac1{|P|}\int_P|\nabla b|^2\right)^{1/2}
 \leq M,
 \end{aligned}
\]
其中可选 \(M>0\) 只依赖 \(d,\lambda,\Lambda\)。注意，所得界控制了 \(Q\) 内每个子立方体，而不只是 \(Q\) 本身；这正是下一步所需的信息。

现在使用\cref{b2:cor:C-local-john-nirenberg}。其假设是初始立方体内的局部可积性，以及所有二进子立方体上的同一平均振荡界。这里 \(b\in H^1(Q)\subset L^1(Q)\)，而刚证得的界对全部轴平行子立方体成立，因而满足该假设。取推论中的维数常数 \(c_d>0\)，令
\(a_0=c_d/(2M)\)，便有
\[
 \frac1{|Q|}\int_Qe^{\pm a_0(b-b_Q)}\leq3.
\]
分别乘回 \(e^{a_0b_Q}\) 和 \(e^{-a_0b_Q}\)，再将两个平均相乘，未知常数因子抵消，得到所列乘积界。附录 C 的局部推论从停止立方体和指数尾的积分直接得到；它的证明不使用 Harnack 比较，因此此处没有循环依赖。
\end{proof}

有了正负平均的连接，只差把负幂平均变成下界。最方便的做法是对倒数使用已证的次解上界，而不是重新进行一遍负幂迭代。

\Needspace{6\baselineskip}
\begin{theorem}[线性正解比较的补充完整证明]
\label{b2:thm:H2-linear-harnack}
设 \(d\geq2\)，\(x_0\in\mathbb R^d\)、\(R>0\)，记 \(B_t=B(x_0,t)\)。可测实矩阵场 \(A\) 在 \(B_{2R}\) 中满足
\(A(x)\xi\cdot\xi\geq\lambda|\xi|^2\)、\(|A(x)\xi|\leq\Lambda|\xi|\)，其中 \(0<\lambda\leq\Lambda<\infty\)，两式对几乎处处的 \(x\) 和全部 \(\xi\in\mathbb R^d\) 成立。若实函数 \(u\in H^1_{\mathrm{loc}}(B_{2R})\) 几乎处处非负，并且
\(\int A\nabla u\cdot\nabla\varphi=0\) 对所有实值 \(\varphi\in C_c^\infty(B_{2R})\) 成立，则
\[
 \operatorname*{ess\,sup}_{B_R}u
 \leq C\operatorname*{ess\,inf}_{B_R}u,
\]
其中 \(C\) 只依赖 \(d,\lambda,\Lambda\)。
\end{theorem}
\begin{proof}
先固定 \(\varepsilon>0\)，令 \(v=u+\varepsilon\)。其倒数 \(w=v^{-1}\) 属于局部 \(H^1\)，而且是非负弱次解。为确认不等式方向，在 \(v\) 的弱方程中使用
\(v^{-2}\cdot\varphi\)，其中 \(0\leq\varphi\in C_c^\infty(B_{2R})\)。这个测试合法，因为 \(v\geq\varepsilon\)。展开导数，得到
\[
 \begin{aligned}
 \int A\nabla w\cdot\nabla\varphi
 &=-\int v^{-2}A\nabla v\cdot\nabla\varphi\\
 &=-2\int v^{-3}\varphi A\nabla v\cdot\nabla v\leq0.
 \end{aligned}
\]
最后一项可积，因为 \(v^{-3}\) 有界、\(\nabla v\) 局部平方可积；其符号只使用椭圆性，不要求矩阵对称。

取 \(4Q\Subset B_{2R}\)。低次平均上界可从球转到立方体：在 \(\overline{Q/2}\) 各点放置半径 \(\ell(Q)/16\) 的球，二倍球仍在 \(Q\) 内，并选有限子覆盖。逐球使用\cref{b2:lem:H2-small-power}，球尺度因子与 \(|Q|^{-1}\) 只相差一个维数常数。因此，对上一个命题给出的同一个 \(a_0>0\)，有
\[
 \begin{aligned}
 \operatorname*{ess\,sup}_{Q/2}v
 &\leq C\left(\frac1{|Q|}\int_Qv^{a_0}\right)^{1/a_0},\\
 \operatorname*{ess\,sup}_{Q/2}v^{-1}
 &\leq C\left(\frac1{|Q|}\int_Qv^{-a_0}\right)^{1/a_0}.
 \end{aligned}
\]
这里 \(v\) 和 \(v^{-1}\) 都是非负次解。相乘后使用对数连接，得到
\[
 \operatorname*{ess\,sup}_{Q/2}v
 \leq H\operatorname*{ess\,inf}_{Q/2}v,
\]
其中 \(H\geq1\) 只依赖 \(d,\lambda,\Lambda\)。左边乘积可以改写成上、下界之比，是因为 \(v\geq\varepsilon\)，且它在 \(Q/2\) 上已有有限上界。

下面把这些小立方体的比较接成整球比较。缩放到 \(R=1\)、\(x_0=0\)，取
\(\ell=1/(10\sqrt d)\)、\(h=\ell/16\)。在 \(\overline{B_1}\) 选一个极大的有限 \(h\)-分离点集 \(c_1,\ldots,c_N\)。极大性说明每个点到某个 \(c_i\) 的距离不超过 \(h\)；而两两不交的球 \(B(c_i,h/2)\) 都在 \(B_{1+h/2}\) 内，所以
\(N\leq(1+2/h)^d\)。令 \(Q_i\) 为中心 \(c_i\)、边长 \(\ell\) 的轴平行立方体。各 \(4Q_i\) 都紧包含于 \(B_2\)，各 \(E_i=Q_i/2\) 覆盖 \(\overline{B_1}\)。

把距离不超过 \(3h\) 的两个中心连成边。所得有限图连通：任取两个中心，沿它们之间位于 \(\overline{B_1}\) 的线段，以不超过 \(h\) 的步长取点，再把每个取样点换成距离不超过 \(h\) 的某个中心；相邻中心的距离便不超过 \(3h\)。每一条边对应的 \(E_i,E_j\) 有正测度重叠，因为
\(3h<\ell/2\)。记
\(S_i=\operatorname*{ess\,sup}_{E_i}v\)、
\(m_i=\operatorname*{ess\,inf}_{E_i}v\)。各片上有 \(S_i\leq Hm_i\)，相邻片的重叠又给出 \(m_i\leq S_j\)。于是沿一条不重复顶点的链，得到
\(S_i\leq H^Nm_j\)。这个论证只比较几乎处处的上下界；有限多个零测例外的并仍为零测集。

选择使 \(S_i\) 最大和使 \(m_j\) 最小的两片，再利用覆盖，便得
\[
 \operatorname*{ess\,sup}_{B_1}v
 \leq\max_iS_i
 \leq H^N\min_jm_j
 \leq H^N\operatorname*{ess\,inf}_{B_1}v.
\]
因为 \(N\) 只受维数控制，常数仍只依赖 \(d,\lambda,\Lambda\)。恢复尺度后，令 \(\varepsilon\downarrow0\)。本性上确界和本性下确界在加常数时都平移同一常数，因此得到 \(u\) 的结论，即使其下确界为零也不例外。
\end{proof}

证明中的加常数步骤避免了预先假设解严格为正。它还显示，正幂上界、倒数次解和对数振荡承担不同任务：前者给局部最大值，倒数把同一个结论转成最小值，对数连接则保证两边平均可以相消。只完成其中的正幂迭代，还不足以得到 Harnack 不等式。

\subsection{De Giorgi--Nash--Moser 理论与 Harnack 不等式}
\label{b2:subsec:H20104}

本节补全了\cref{b2:thm:16-measurable-harnack}中的线性比较路线。相关公开论文 \cite[Theorems 1.1--1.3 and Section 2]{Li2017Harnack}讨论了局部上界、正性传播与 Harnack 比较之间的关系；这里实际使用的指数尾则由\appref{b2:app:C}在书内证明。上述论证始终保留了实标量、非负解、内球和一致椭圆假设。

\ProseParagraphBreak

Harnack 比较对有符号解的作用，要通过减去常数来显现。在外球上记 \(M=\operatorname*{ess\,sup}u\)、\(m=\operatorname*{ess\,inf}u\)，则 \(u-m\) 和 \(M-u\) 是两个非负弱解。分别应用比较并相加，便得到较小球上的严格振荡收缩；\subsecref{b2:subsec:160404}已经给出收缩因子、二进半径迭代以及从本性上下界构造 Hölder 连续代表的完整步骤。现在该推导所需的比较也已在本节证明，无须再把连续性当作额外输入。

\ProseParagraphBreak

正解比较还有一个可以直接读出的后果：非负解若在某个小球上具有零的本性下界，就必须在该球上消失；连通性将这种消失传到整个区域。这是标量方程中强最大值原理的一种内部形式。

\Needspace{6\baselineskip}
\begin{corollary}[连通区域中的正性与强最大值原理]
\label{b2:cor:H2-connected-positivity}
设 \(d\geq2\)，\(\Omega\subset\mathbb R^d\) 为连通开集，\(A:\Omega\to\mathbb R^{d\times d}\) 为可测实矩阵场，并满足
\(A(x)\xi\cdot\xi\geq\lambda|\xi|^2\)、\(|A(x)\xi|\leq\Lambda|\xi|\) 对几乎处处的 \(x\) 和全部实 \(\xi\) 成立，其中 \(0<\lambda\leq\Lambda<\infty\)。若实函数 \(u\in H^1_{\mathrm{loc}}(\Omega)\) 几乎处处非负，且
\(\int A\nabla u\cdot\nabla\varphi=0\) 对全部实值 \(\varphi\in C_c^\infty(\Omega)\) 成立，则或者 \(u=0\) 几乎处处，或者对每个 \(B_{2r}(x)\Subset\Omega\)，都有
\[
 \operatorname*{ess\,inf}_{B_r(x)}u>0.
\]
因而，非零解的局部 Hölder 连续代表在 \(\Omega\) 中处处为正。更一般地，若实弱解 \(z\in H^1_{\mathrm{loc}}(\Omega)\) 满足 \(z\leq M\) 几乎处处，其中 \(M\in\mathbb R\)，且其连续代表在某个内点取值 \(M\)，则 \(z=M\) 几乎处处，其连续代表处处等于 \(M\)。
\end{corollary}
\begin{proof}
令 \(Z\) 为所有具有一个邻域、使 \(u\) 在该邻域上几乎处处为零的点。按定义，\(Z\) 是相对开集。它也相对闭：若 \(x\in\overline Z\cap\Omega\)，取 \(B_{2r}(x)\Subset\Omega\)。球 \(B_r(x)\) 与 \(Z\) 的交集含有非空开集，所以 \(u\) 在一个正测度子集上为零；由非负性，\(\operatorname*{ess\,inf}_{B_r(x)}u=0\)。Harnack 比较给出 \(u=0\) 几乎处处于 \(B_r(x)\)，故 \(x\in Z\)。连通性说明 \(Z\) 为空或等于 \(\Omega\)；后一情形用可数子覆盖合并零测例外，得到 \(u=0\) 几乎处处于整个区域。

若 \(u\) 不是零函数，\(Z\) 必为空。任意 \(B_{2r}(x)\Subset\Omega\) 的内球若有零的本性下界，Harnack 比较会使该内球成为 \(Z\) 的一部分，矛盾。因此本性下界严格为正。连续代表与原函数几乎处处相等，故在内球上也不能低于这一正下界；否则连续性会给出一个正测度的违反集合。

最后对 \(u=M-z\) 应用这个二择结论。它是非负弱解，其连续代表在指定内点为零，因而不可能属于处处严格为正的情形。于是 \(z=M\) 几乎处处；连续性再把等式延伸到每一点。
\end{proof}

上述局部有界、正解比较和 Hölder 正则性通常合称 \term[DeGiorgiNashMoserLilun]{De Giorgi--Nash--Moser 理论}[De Giorgi--Nash--Moser theory]。名字相连并不意味着三个结论可以互相省略。对于椭圆系统、复值方程或一般非齐次方程，标量正部测试、序结构或齐次缩放可能失效。下一节转向 \(p\) 次增长方程时，会重新建立相应的截断估计和局部上界，而不把本节的线性 Harnack 定理直接套入。
